\section{The graph and completeness}\label{graph:section}

Fix a source instance and the construction parameters from
\Cref{sec:source}. Each vertex will specify source constraints on the
pairs of \(d\) positions and weights on the labels available at each
position. We first define when local labels fit together and use their
coordinate sums to define a norm. A satisfying source labeling will
then provide a linear functional whose values above \(t\) select an
independent set.

\subsection{Queries and compatible labels}

Let $\mathcal E=\binom{[d]}2$ be the set of unordered pairs of distinct
positions.  A query of type $j\in[d]$ has the form
\[
 i=\bigl(j,(q_{jk})_{k\in[d]\setminus\{j\}}\bigr),
 \qquad
 q_{jk}\in
 \begin{cases}
  V,&k<j,\\
  U,&k>j.
 \end{cases}
\]
The tuple is ordered by increasing $k$.  Let $\mathcal Q$ be the set
of all such queries of all types, and let $P(i)\subseteq U\cup V$ be
the set of distinct variables appearing in $i$.  The side of each
variable is part of its identity, since $U$ and $V$ are disjoint.
Repeated entries in the tuple contribute only one variable to $P(i)$;
thus $|P(i)|\leq d-1$.

A label $\lambda$ for $i$ assigns a symbol from the appropriate
alphabet to each variable of $P(i)$ and satisfies every source
constraint occurrence whose two endpoints belong to $P(i)$.  Write
$\Lambda_i$ for the set of these labels.  This set may be empty, for
example when constraints internal to the scope are inconsistent.
There are at most $q^{|P(i)|}\leq r$ labels.  Fix orders on variables
and alphabet symbols, and list the valid assignments in lexicographic
order to define an injection
\[
 \operatorname{slot}_i:\Lambda_i\longrightarrow[r].
\]
This is the empty map when $\Lambda_i=\varnothing$.  Slots outside
its image will play no role in an increment.

Define the finite coordinate set
\[
 \mathcal P=\{(i,\lambda):i\in\mathcal Q,\ \lambda\in\Lambda_i\}.
\]
A \emph{compatible selection} is a subset $I\subseteq\mathcal P$
with at most one label per query, such that the selected labels agree
on overlaps and their induced labeling satisfies every source
constraint whose two endpoints lie in
\[
 P(I)=\bigcup_{(i,\lambda)\in I}P(i).
\]
Let $\mathcal I$ be the family of compatible selections, including
the empty selection. Agreement on overlaps alone is not enough:
constraints joining variables from different selected scopes must
also be satisfied. Thus labels on any two selected queries satisfy every source
constraint between their scopes.
Compatibility imposes no condition on a constraint with an endpoint
outside $P(I)$; the induced labeling need not extend to a satisfying
labeling of the entire instance.

For $z\in\R^{\mathcal P}$, define
\begin{equation}\label{graph:norm-definition}
 \norm{z}=\max_{I\in\mathcal I}
       \left|\sum_{(i,\lambda)\in I}z_{i,\lambda}\right|.
\end{equation}
For a query $i$ and a vector $s\in\R^r$, let $z(i,s)$ be supported on
its query block, with
\[
 z(i,s)_{i,\lambda}=s_{\operatorname{slot}_i(\lambda)}
 \quad\text{for }\lambda\in\Lambda_i.
\]
All coordinates in other query blocks are zero.  In particular,
$z(i,s)=0$ for a query with no valid label.

\begin{lemma}[Compatibility norm]\label{graph:norm}
The family $\mathcal I$ is hereditary and contains every coordinate
singleton.  Formula~\eqref{graph:norm-definition} defines a norm.
For each fixed query $i$, the increment is linear in $s$ and obeys
\begin{equation}\label{graph:increment}
 \norm{z(i,s)}\leq\|s\|_\infty.
\end{equation}
\end{lemma}
\begin{proof}
Removing labels preserves agreement and can only shrink the scope
union.  Every constraint internal to the smaller union was therefore
satisfied before removal.  This proves heredity.  Every singleton
belongs to $\mathcal I$ by the definition of a valid query label.
The maximum of absolute linear forms is absolutely homogeneous and
satisfies the triangle inequality; the singleton forms ensure that
it vanishes only at zero.  If the coordinate set is empty, this is
the unique norm on the zero-dimensional space.  Finally, a compatible
selection contains at most one coordinate of the $i$-block, proving
\eqref{graph:increment}.  Linearity follows from the coordinate formula.
\end{proof}

\subsection{Vertices and edges}

A seed tuple is an indexed family
$\mathbf c=(c_e)_{e\in\mathcal E}\in[M]^{\mathcal E}$.
For $e=\{j,k\}$ with $j<k$, seed $c_e$ gives endpoint $x_{c_e}$ to
position $j$ and endpoint $y_{c_e}$ to position $k$.  It thereby
specifies queries $i_j(\mathbf c)$ through
\[
 q_{jk}=x_{c_{\{j,k\}}},\qquad
 q_{kj}=y_{c_{\{j,k\}}}\qquad(j<k).
\]
A query retains its type and its ordered endpoint questions.  It does
not retain the sampled occurrence indices, the sampled projection
maps, or the opposite endpoint in any pair.  Different seed tuples
can give the same queries.  Queries of different position types are
nevertheless distinct, even if some endpoint values coincide.
The source instance used to define labels and slots is fixed; the
seed-dependent data in a query are only its endpoint questions.

Use the symmetric grid
\begin{equation}\label{graph:grid}
 \mathcal T=\left\{\frac{2u}{p}:u=-d,-d+1,\ldots,d\right\},
 \qquad p=2d+1.
\end{equation}
These are the midpoints of the $p$ equal consecutive subintervals of
$[-1,1]$.  The vertex set consists of all pairs
\[
 v=(\mathbf c,\mathbf s),\qquad
 \mathbf c\in[M]^{\mathcal E},\quad
 \mathbf s=(s_1,\ldots,s_d)\in(\mathcal T^r)^d.
\]
Each indexed pair is one vertex, including when another pair produces
the same queries or the same vector
\begin{equation}\label{graph:sum-vector}
 S_v=\sum_{j=1}^d z(i_j(\mathbf c),s_j).
\end{equation}
For two distinct vertices, declare
\begin{equation}\label{graph:edge}
 \{v,w\}\in E(G)
 \quad\Longleftrightarrow\quad
 \norm{S_v+S_w}\leq2t.
\end{equation}
This defines a finite simple undirected unweighted graph.  Its order is
\begin{equation}\label{graph:size}
 n=|V(G)|=M^{\binom d2}p^{rd}\geq p>2m,
\end{equation}
where the inequalities use $M,r\geq1$ and
\Cref{src:dimension-bounds,src:construction-parameters}. Choosing a uniform vertex is exactly
choosing all seeds independently and uniformly from $[M]$ and all
weight entries independently and uniformly from $\mathcal T$.
There is no identification of equal vectors in this probability law.

\begin{lemma}[Local evaluation of the norm]\label{graph:local-norm}
Let $\mathcal B\subseteq\mathcal Q$ be a set of query blocks
supporting $z\in\R^{\mathcal P}$.  Then
\begin{equation}\label{graph:local-formula}
 \norm{z}=\max_{\substack{I\in\mathcal I:\\
                          (i,\lambda)\in I\ \Longrightarrow\ i\in\mathcal B}}
       \left|\sum_{(i,\lambda)\in I}z_{i,\lambda}\right|.
\end{equation}
For $z=S_v+S_w$, this maximum can be evaluated by at most
$(1+r)^{2d}$ choices, each checked by a scan of the source constraints.
\end{lemma}
\begin{proof}
Intersect any compatible selection with the blocks in $\mathcal B$.
Heredity preserves compatibility, and the sum is unchanged because
all other blocks of $z$ vanish.  Conversely, each selection on the
right of \eqref{graph:local-formula} is a selection in the full
maximum.  This proves equality.

For a pair of vertices there are at most $2d$ distinct query blocks.
First aggregate the coordinate contributions of all occurrences of
each repeated query.  For each distinct query $i$, enumerate the at
most $q^{d-1}\leq r$ assignments on its scope, and keep exactly those
satisfying every source constraint internal to that scope.  This
computes $\Lambda_i$ and its slots.  Next choose either no label or
one of these labels in each block.  There are at most $(1+r)^{2d}$
choices.  For each choice, test agreement on repeated variables and
scan the source list to test every constraint internal to the union
of the chosen scopes.  This union has at most $2d(d-1)$ variables.
The surviving choices are precisely the selections on the right of
\eqref{graph:local-formula}, and their absolute coordinate sums give
the required maximum.  Neither the labels nor the compatibility test
requires satisfying any constraint outside the indicated scope union.
\end{proof}

\begin{proposition}[Deterministic polynomial construction]
\label{graph:polynomial}
For a fixed integer $m\geq4$, a fixed even square $d$ satisfying
\Cref{src:dimension-bounds}, and fixed source alphabets,
the graph $G$ can be output explicitly in deterministic time
polynomial in the source bit length.  Every edge test is exact.
\end{proposition}
\begin{proof}
Let $L$ be the source bit length. The fixed integers $m,d$ determine
$p$ and the rational number $t$ by \Cref{src:construction-parameters}.
The fixed alphabets determine $q$ and $r=q^{d-1}$. All these
quantities are therefore computable constants independent of $L$.
Since the
constraint list is explicit, $M=O(L)$, and
\eqref{graph:size} gives $n=O(L^{\binom d2})$.  A vertex description
contains a fixed number of source indices and grid entries.  Thus all
vertices can be enumerated and assigned identifiers of
$O(1+\log n)$ bits.

Apply \Cref{graph:local-norm} to each unordered pair of distinct
vertices.  All label and selection enumerations have a fixed number
of choices.  Variable comparisons, sorting, constraint scans, and
projection-map evaluations take polynomial bit time in $L$.
The arithmetic of the edge test uses only rational grid entries.
In fact, every coordinate of $S_v+S_w$ has denominator $p$, after
summing all repeated-block contributions.  The local maximum is
therefore $\zeta/p$ for a nonnegative integer $\zeta$.  By
\eqref{graph:increment} and the triangle inequality,
$\zeta/p\leq2d$, so its numerator and denominator have bit lengths
bounded solely in terms of the fixed parameters.  Since
$t=\sqrt d/(4m)$ and $\sqrt d$ is an integer, the edge test is precisely
\[
 2m\zeta\leq p\sqrt d.
\]
Both sides are integers, so the test involves no approximation of a
real number.

Checking all pairs and writing every qualifying edge takes
\[
 O\bigl((1+n^2)\operatorname{poly}(L+\log n)\bigr)
\]
bit operations, with constants depending only on the fixed parameters
and alphabets.  Formula~\eqref{graph:size} also computes $n$ by integer
arithmetic with $O(1+\log n)$ output bits.  Thus the construction has
polynomial bit complexity, including its explicit vertex and edge
lists.  It runs on every source instance with these alphabets,
regardless of whether the gap promise holds.
\end{proof}

\subsection{Completeness}

A satisfying source assignment will give a linear functional whose
value exceeds \(t\) on an independent set. To show that this set
occupies nearly half the vertices, we first bound the corresponding
tail of a sum of independent uniform grid variables.

\begin{lemma}[Midpoint-grid tail]\label{graph:anticoncentration}
If $T_1,\ldots,T_d$ are independent uniform elements of $\mathcal T$,
then
\begin{equation}\label{graph:tail}
 \Pr\left[\sum_{j=1}^dT_j>t\right]
 >\frac12-\sqrt2\left(a+\frac1{\sqrt d}\right)
 >\frac12-\frac1m.
\end{equation}
\end{lemma}
\begin{proof}
Take independent uniform $C_1,\ldots,C_d$ on $[-1,1]$, and round each
to the midpoint of its subinterval in the partition defining
$\mathcal T$.  Boundary choices have probability zero.  The rounded
variables have the required independent laws, and
\[
 \left|\sum_jT_j-\sum_jC_j\right|\leq\frac dp<\frac12.
\]
Put $A=t+d/p=a\sqrt d+d/p$.  We bound the central probability
$\Pr[|\sum_jC_j|\leq A]$.

Each $C_j$ may be generated as $(L_j+U_j)/2$, where $L_j$ is a fair
sign and $U_j$ is uniform on $[-1,1]$, all independent.  The two sign
choices give uniform distributions on the two half intervals of
$[-1,1]$, so this is indeed the uniform distribution on the full
interval.  Conditional on $(U_1,\ldots,U_d)$, the event
$|\sum_jC_j|\leq A$ places $\sum_jL_j$ in an interval of length
\[
 4A=4a\sqrt d+4d/p<4a\sqrt d+2.
\]
The possible sign sums are spaced by two.  An interval of length
$D$ contains at most $D/2+1$ values from such a lattice, so this
interval contains fewer than $2a\sqrt d+2$ possible sign sums.

Write $d=2d_0$.  The sign-sum probabilities are
$4^{-d_0}\binom{2d_0}{v}$, for $0\leq v\leq2d_0$.
Consecutive ratios $(2d_0-v)/(v+1)$ show that the largest is the
central one.  It satisfies
\[
 4^{-d_0}\binom{2d_0}{d_0}
 =\prod_{u=1}^{d_0}\left(1-\frac1{2u}\right)
 \leq\frac1{\sqrt{d_0+1}}
 <\frac{\sqrt2}{\sqrt d}.
\]
For the middle inequality, square the product and use
\[
 \left(1-\frac1{2u}\right)^2\leq\frac{u}{u+1}
 \qquad(u\geq1);
\]
indeed the right side minus the left side is
$(3u-1)/(4u^2(u+1))>0$, and the resulting product telescopes.
The conditional central probability, and therefore its average,
is bounded by
\[
 \Pr\left[\left|\sum_jC_j\right|\leq A\right]
 <\sqrt2\left(2a+\frac2{\sqrt d}\right).
\]
Symmetry gives
\[
 \Pr\left[\sum_jC_j>A\right]
 =\frac12\left(1-
    \Pr\left[\left|\sum_jC_j\right|\leq A\right]\right)
 >\frac12-\sqrt2\left(a+\frac1{\sqrt d}\right).
\]
On this event, rounding leaves $\sum_jT_j>t$.  Finally,
$\sqrt d>4m=1/a$ by \Cref{src:dimension-bounds}, and hence
\[
 \sqrt2\left(a+\frac1{\sqrt d}\right)
 <2\sqrt2\,a<3a=\frac3{4m}<\frac1m.
\]
This proves both inequalities in \eqref{graph:tail}.
\end{proof}

\begin{proposition}[Completeness]\label{graph:completeness}
If $\val(\Phi)=1$, then $G$ has an independent set of size greater
than $(1/2-1/m)n$.  Consequently,
\[
 \tau(G)<\left(\frac12+\frac1m\right)n.
\]
\end{proposition}
\begin{proof}
Let $A$ be a labeling satisfying every source constraint.  Its
restriction $\lambda_{A,i}$ to $P(i)$ belongs to $\Lambda_i$ for
every query $i$.  The selection
\[
 I_A=\{(i,\lambda_{A,i}):i\in\mathcal Q\}
\]
belongs to $\mathcal I$: the labels agree on overlaps, and every
constraint internal to their scope union is satisfied by $A$.
Keep precisely the vertices for which
\begin{equation}\label{graph:keep-rule}
 \sum_{j=1}^d
 s_{j,\operatorname{slot}_{i_j(\mathbf c)}
                  (\lambda_{A,i_j(\mathbf c)})}>t.
\end{equation}
For a uniform vertex, conditional on the seed tuple, the selected
slot in each weight vector is fixed.  Its entries are independent
uniform elements of $\mathcal T$, one from each independent $s_j$.
Coincident endpoint questions or slot numbers do not change this
independence.  By \Cref{graph:anticoncentration}, more than a
$(1/2-1/m)$ fraction of vertices are kept.

For two distinct kept vertices $v,w$, evaluating $S_v+S_w$ on $I_A$
gives exactly the sum of their two expressions in
\eqref{graph:keep-rule}.  This includes all contributions even when
the same query appears in both vertices.  The sum exceeds $2t$, so
\eqref{graph:norm-definition} rules out the edge in
\eqref{graph:edge}.  The kept vertices are therefore independent.
Their complement is a vertex cover of size less than
$(1/2+1/m)n$.
\end{proof}
